%! Measuring Center and Spread — Unit 2, Lesson 2.2 — Homework
% This is the lesson's SCORED individual practice and the LAST component in the packet.
% Unit 2 timing (5 / 45 / 5 / 5): ASSIGNED at Close & Assign and DONE OUTSIDE CLASS — not
% started in class. Due at the start of the first class after two study halls; 2 pts per item.
% THIRD CONTEXT: the notes teach on Maple Grove's bus routes, the Guided Practice runs on two
% towns' temperatures, and these items are on Oak Street Pizza's delivery times.
% TWO PAGES MAX: two boxes — items 1-3 on page 1, items 4-6 and the spiralbox on page 2 — so no
% item breaks across a page.
% ITEMS: 1 center, interpreted (x-bar with Sigma, median) · 2 SPIRAL to Lesson 1.2 (statistic
% vs parameter) · 3 five-number summary, range, IQR, interpreted · 4 CORE PROCEDURE: s from a
% deviations table · 5 the CONTRAST PAIR: (a) same mean, different s; (b) different mean, same
% s — the pair that surfaces "bigger values, bigger SD" · 6 the CRUX HEAD-ON + JUSTIFY: every
% delivery +8 min; Jordan says s went up.
% THE DATA — Oak Street Pizza, six Friday deliveries (min): 26, 26, 28, 32, 32, 36
%   sum 180, n 6, mean 30; median (28+32)/2 = 30; lower half 26,26,28 -> Q1 = 26;
%   upper half 32,32,36 -> Q3 = 32; IQR 6; range 36-26 = 10
%   deviations -4,-4,-2,2,2,6 (sum 0); squares 16,16,4,4,4,36 (sum 80); s^2 = 80/5 = 16; s = 4
%   Main Street: mean 30, s 11. Hilltop: mean 45, s 4. Snowy Friday (+8): mean 38, s 4.
% PAGE LOCKSTEP: the key differs only in -key, the header, \blank -> \ans, and each
% \par\writelines{n} -> n \ansline; every work block is byte-identical.
\documentclass[12pt]{article}
\usepackage{saar-article}
\usepackage{saar-boxes}
\newcommand{\TallMath}[1]{$\displaystyle #1\rule[-1.4em]{0pt}{3.2em}$}
% Word bank strip for every fill-in cluster. Defined per document; the key inherits it.
\newcommand{\wordbank}[1]{\par\vspace{2pt}\noindent\fcolorbox{goldacc}{goldbg}{\parbox{\dimexpr\linewidth-2\fboxsep-2\fboxrule\relax}{\small\textbf{Word bank:} #1}}\par\vspace{4pt}}

\begin{document}

\pageheader{Unit 2, Lesson 2.2}{Homework}
% NAMESTRIP: no \namedateperiod here. The student writes their name once, on the cover.

\vspace{0.04in}

% Authored identically in homework/ and homework_key/.
\begin{remindbox}
\small
\textbf{This is your graded homework.} It is scored out of 12 --- 2 points per item --- and due
at the start of the first class after two study halls. It is done on your own, outside class;
the \emph{Keep in Mind} box on the cover is the page to flip back to.
\end{remindbox}

\vspace{0.03in}

\boxguard[20]
\begin{scenariobox}[Oak Street Pizza]{cerulean}
Oak Street Pizza advertises delivery ``in about half an hour.'' One Friday night the manager
timed six deliveries, in minutes:

\vspace{2pt}
\centerline{$26,\quad 26,\quad 28,\quad 32,\quad 32,\quad 36$}
\end{scenariobox}

\vspace{0.03in}

\boxguard
\begin{notesbox}{Center and Spread}
Every answer says what the number means for Oak Street's deliveries.
\wordbank{statistic \;\textbullet\; parameter \;\textbullet\; the 6 timed Friday deliveries \;\textbullet\; all of Oak Street's Friday deliveries \;\textbullet\; $6$ \;\textbullet\; $10$ \;\textbullet\; $26$ \;\textbullet\; $30$ \;\textbullet\; $32$ \;\textbullet\; $36$}
\begin{enumerate}[leftmargin=1.6em, itemsep=5pt, topsep=2pt]

  \item Find the mean delivery time. Write it with $\bar{x}$ and $\Sigma$.

        \begin{work}
          \bar{x} = \dfrac{\sum x}{n} = \dfrac{180}{6} = 30 \text{ minutes}
        \end{work}

        Median: \blank{3.4cm}
        \par\vspace{6pt}
        What does $\bar{x}$ tell a customer about Oak Street's deliveries?
        \par\writeline

  \item \emph{(Lesson 1.2.)} These six deliveries are a sample of all of Oak Street's Friday
        deliveries. Is $\bar{x} = 30$ minutes a parameter or a statistic, and who does it
        describe?

        \par\vspace{4pt}
        Kind: \blank{2.6cm} \hfill Describes: \blank{7cm}

  \item Complete the five-number summary (use the median of each half), then find the range
        and the IQR.

        \begin{center}
        \renewcommand{\arraystretch}{1.2}
        \begin{tabular}{@{} C{2.2cm} C{2.2cm} C{2.2cm} C{2.2cm} C{2.2cm} @{}}
        \toprule
        \textbf{Min} & $\boldsymbol{Q_1}$ & \textbf{Median} & $\boldsymbol{Q_3}$ & \textbf{Max} \\
        \midrule
        \blank{1.4cm} & \blank{1.4cm} & \blank{1.4cm} & \blank{1.4cm} & \blank{1.4cm} \\
        \bottomrule
        \end{tabular}
        \end{center}

        Range: \blank{3.2cm} \qquad IQR: \blank{3.2cm}
        \par\vspace{6pt}
        Interpret the IQR for Oak Street's customers.
        \par\writeline

\end{enumerate}
\end{notesbox}

\boxguard[30]
\begin{notesbox}{Center and Spread, continued}
\wordbank{$-4$ \;\textbullet\; $-2$ \;\textbullet\; $0$ \;\textbullet\; $2$ \;\textbullet\; $4$ \;\textbullet\; $6$ \;\textbullet\; $16$ \;\textbullet\; $36$ \;\textbullet\; $38$ \;\textbullet\; $80$}
\begin{enumerate}[leftmargin=1.6em, itemsep=5pt, topsep=2pt, start=4]

  \item Complete the deviations table ($\bar{x} = 30$), then find the variance and the
        standard deviation. Interpret $s$.

        \begin{center}
        \renewcommand{\arraystretch}{1.35}
        \begin{tabular}{@{} l C{1.2cm} C{1.2cm} C{1.2cm} C{1.2cm} C{1.2cm} C{1.2cm} C{1.3cm} @{}}
        \toprule
        \textbf{Time} $x$ & $26$ & $26$ & $28$ & $32$ & $32$ & $36$ & Sum \\
        \midrule
        $x - \bar{x}$ & \blank{0.9cm} & \blank{0.9cm} & \blank{0.9cm} & \blank{0.9cm} & \blank{0.9cm} & \blank{0.9cm} & \blank{0.9cm} \\
        $(x - \bar{x})^2$ & \blank{0.9cm} & \blank{0.9cm} & \blank{0.9cm} & \blank{0.9cm} & \blank{0.9cm} & \blank{0.9cm} & \blank{0.9cm} \\
        \bottomrule
        \end{tabular}
        \end{center}

        Variance $s^2 =$ \blank{3.2cm} \qquad Standard deviation $s =$ \blank{3.2cm}
        \par\vspace{6pt}
        Interpret $s$ in context.
        \par\writeline

  \item Two other shops timed Friday deliveries too. \textbf{Main Street Pizza:}
        $\bar{x} = 30$ min, $s = 11$ min. \textbf{Hilltop Pizza:} $\bar{x} = 45$ min,
        $s = 4$ min.

        \par\vspace{3pt}
        \textbf{(a)} Oak Street and Main Street have the same mean. Whose delivery times are
        more predictable? Explain.
        \par\writeline

        \textbf{(b)} Hilltop's deliveries take 15 minutes longer on average. Are Hilltop's
        times more spread out than Oak Street's? Explain.
        \par\writeline

  \item On a snowy Friday every Oak Street delivery took exactly $8$ minutes longer. Jordan
        says, ``Bigger times, so they're more spread out --- $s$ went up.'' Give the new
        $\bar{x}$ and $s$, and explain what Jordan has mixed up.

        \par\vspace{4pt}
        New $\bar{x}$: \blank{3cm} \qquad New $s$: \blank{3cm}
        \par\writeline\par\writeline

\end{enumerate}
\end{notesbox}

\boxguard[5]   % the short spiralbox needs about five baselines; a larger guard can push it to a third page
\begin{spiralbox}
\textbf{Next lesson: Boxplots, Outliers, and Resistance.} Item~3's five-number summary is what a
boxplot is drawn from; next class we draw one and test for outliers.
\end{spiralbox}

\end{document}
